
\begin{table}[h]
\caption{An example of how a reasoning path is augmented into four formats of training data with different prompts (in input) and answer styles (in output). Specifically, the \textit{CoT prompting examples} used for each tasks are listed in Appendix~\ref{sec:app_prompt}. The \textit{Standard prompting examples} are the same question-answer pairs with \textit{CoT prompting examples}, except that reasoning is removed.}
% \vspace{1.5em}
\centering
% \small
\resizebox{\textwidth}{!}{
\begin{tabular}{l}
\toprule
\textbf{Question}: Amy is 10 years old. Jake is 8 years old. Alex's age is right in the middle. How old is Alex?\\
\textbf{Selected Chain-of-Thought}: Amy is 10 years old. Jake is 8 years old. Alex's age is in the middle of\\ Amy and Jake, so Alex is ( 8 + 10 ) / 2 = 9 years old. The answer is 9.\\
\midrule
\textbf{Mixed-formats of training data}:\\
\textbf{Format 1}: \textbf{Input}: \textit{[CoT prompting examples]} + `\textbackslash n' + \textit{[Question]} +  `\textbackslash n' + `A:'\\
\textbf{Output}: Amy is 10 years old. Jake is 8 years old. Alex's age is in the middle of Amy and Jake, so Alex\\ is ( 8 + 10 ) / 2 = 9 years old. The answer is 9.\\
\\
\textbf{Format 2}: \textbf{Input}: \textit{[Standard prompting examples]} + `\textbackslash n' + \textit{[Question]} +  `\textbackslash n' + `A:'\\
\textbf{Output}: The answer is 9.\\
\\
\textbf{Format 3}: \textbf{Input}: \textit{[Question]} +  `\textbackslash n' + `A: Let's think step by step.'\\
\textbf{Output}: Amy is 10 years old. Jake is 8 years old. Alex's age is in the middle of Amy and Jake, so Alex\\ is ( 8 + 10 ) / 2 = 9 years old. The answer is 9.\\
\\
\textbf{Format 4}: \textbf{Input}: \textit{[Question]} +  `\textbackslash n' + `A:'\\
\textbf{Output}: The answer is 9.\\
\bottomrule
\end{tabular}
}
\label{tab:mix_format}
% \vspace{-1em}
\end{table}
